\section{实验环境}

\begin{itemize}
    \item 硬件环境 1（本地测试与 QEMU 宿主环境）
    \begin{itemize}
        \item 处理器：13th Gen Intel(R) Core(TM) i5-13420H
        \item 内存：16 GB
    \end{itemize}

    \item 硬件环境 2
    \begin{itemize}
        \item 处理器：Apple M4
        \item 内存：16 GB
    \end{itemize}

    \item 软件环境 1（本实验正文主环境）
    \begin{itemize}
        \item 操作系统：Windows 11
        \item WSL2：Arch Linux x86\_64
        \item 交叉编译工具链：\verb|riscv64-unknown-linux-gnu-gcc|
        \item 模拟环境：QEMU System RISC-V + OpenSBI + Linux
    \end{itemize}

    \item 软件环境 2
    \begin{itemize}
        \item 操作系统：macOS
        \item Docker Ubuntu 22.04 LTS
    \end{itemize}
\end{itemize}

其中x86主机环境主要用于代码调试、参考性能观察和新增实现的快速验证。
